As a concrete illustration of the general formalism, consider a single
step of the separable walk on a $3 \times 3$ torus with $\rho_x =
\rho_y = 2/3$ and $\delta_x = \delta_y = 0$. Take the walker initially
localised at the origin with coin state $|RU\rangle = |R\rangle \otimes
|U\rangle$ --- that is, moving in the positive direction along both axes.

The coin operator acts first. Each single-axis coin is
\begin{equation}
C\!\left(\tfrac{2}{3}, 0\right)
= \begin{pmatrix}
\sqrt{2/3} & \sqrt{1/3} \\
\sqrt{1/3} & -\sqrt{2/3}
\end{pmatrix},
\label{eq:coin-2-3}
\end{equation}
acting on the basis $\{|\text{negative}\rangle, |\text{positive}\rangle\}$
along each axis. Acting on the initial coin state:
\begin{align}
C_x |R\rangle &= \sqrt{1/3}\,|L\rangle - \sqrt{2/3}\,|R\rangle, \\
C_y |U\rangle &= \sqrt{1/3}\,|D\rangle - \sqrt{2/3}\,|U\rangle.
\end{align}
Taking the tensor product $C_x \otimes C_y$ gives the post-coin state
\begin{equation}
(C_x \otimes C_y)|RU\rangle
= \tfrac{1}{3}|LD\rangle - \tfrac{\sqrt{2}}{3}|LU\rangle
- \tfrac{\sqrt{2}}{3}|RD\rangle + \tfrac{2}{3}|RU\rangle.
\label{eq:post-coin}
\end{equation}

The conditional shift $S = S_x S_y$ then moves each branch by one site
along each axis, in the direction indicated by the coin state. On the
$3 \times 3$ torus, the four destinations are
\begin{align}
|LD\rangle &: (x,y) \to (2, 2), \\
|LU\rangle &: (x,y) \to (2, 1), \\
|RD\rangle &: (x,y) \to (1, 2), \\
|RU\rangle &: (x,y) \to (1, 1).
\end{align}
The state after one full step is therefore
\begin{equation}
|\psi_1\rangle
= \tfrac{1}{3}|2,2\rangle|LD\rangle
- \tfrac{\sqrt{2}}{3}|2,1\rangle|LU\rangle
- \tfrac{\sqrt{2}}{3}|1,2\rangle|RD\rangle
+ \tfrac{2}{3}|1,1\rangle|RU\rangle,
\label{eq:state-after-one-step}
\end{equation}
in which position and coin are now entangled. The norm is preserved:
\begin{equation}
\langle\psi_1|\psi_1\rangle
= \tfrac{1}{9} + \tfrac{2}{9} + \tfrac{2}{9} + \tfrac{4}{9} = 1.
\end{equation}
Figure~\ref{fig:one-step} shows the coin splitting the amplitude over
the four directions (panel a) and the shift distributing the branches
to their respective nodes on the torus (panel b). Unlike the
$\rho = 1/2$ Hadamard case, the split is asymmetric: the
$|RU\rangle$ branch retains amplitude $2/3$, while the two off-diagonal
branches carry $-\sqrt{2}/3$ and the opposite-corner branch carries
$1/3$.
